% Related work, Version B structure, three flows. Pick one. Citations: related_works.bib
% Thesis: the container wall is the sensor. Sound shakes the wall, the object inside changes the shaking,
% we read it remotely at many points and learn to turn it into an image. So: any wall (incl. metal),
% no contact, object at rest, spatial output, cheap and safe hardware.

% =====================================================================
% FLOW 1: simple funnel (background -> limitation -> ours, in every subsection)
% =====================================================================
\section{Related Work}
\label{sec:related}

\subsection{Inferring contents via active sensing}
Seeing inside a closed container usually means sending a wave through its walls.
X-rays are the standard tool in security screening, but they are ionizing, need shielding, and are blocked by dense metal.
Radio waves are safe and pass through walls and packaging.
WiFi can detect people behind a wall~\cite{adib2013see}, and networks trained with a camera as the teacher recover their pose and body shape from radio alone~\cite{zhao2018rfpose,wang2019personinwifi,geng2022densepose}.
Millimeter-wave radar goes further and images objects inside envelopes and boxes~\cite{ganse2021envelopes,dodds2025mmnorm,lu2026geraf,dodds2025waveformer}.
However, most objects reflect radio waves like a mirror, so the radar sees only the surfaces that face it, and the rest must be filled in by surface-normal estimation~\cite{dodds2025mmnorm}, neural fields~\cite{lu2026geraf}, or learned shape completion~\cite{dodds2025waveformer}.
Radio also cannot pass through metal.
Acoustic echoes~\cite{lindell2019acoustic,christensen2020batvision} and optical time-of-flight~\cite{velten2012recovering} can recover hidden geometry too, but they need an open path to the object and their resolution is limited by the wavelength.

Our method sends nothing through the wall.
We vibrate the wall with sound and measure how the object inside changes that vibration.
Because the wall itself carries the signal, the method works for any wall material, including metal, and needs only a speaker, a laser, and a camera.

\subsection{Inferring contents via contact vibration sensing}
Vibration is a rich source of information about objects.
The sound of a struck object reveals its shape and material~\cite{zhang2017shape,gao2023objectfolder,clarke2023realimpact}.
Contact microphones let robots recognize what they hold~\cite{liu2024sonicsense} and where they are touched~\cite{zaielamri2026vibrosense}, and let ordinary surfaces act as touch sensors~\cite{liu2017vibsense,iravantchi2023sawsense}.
In structural engineering, shifts in a structure's frequency response reveal added masses and damage~\cite{dackermann2013identification,ostachowicz2002location,azuara2021damage}.
These methods attach sensors to the object and predict a label, a property, or a location.
The Boombox~\cite{chen2021boombox} comes closest to our goal.
Contact microphones on a box record objects as they are dropped in, and a network predicts an image of where they land.
Its signal comes from the impact, so it only works while the object is moving.
We excite the container with a speaker instead, which lets us image objects at rest, and we sense the wall from a distance, so nothing touches the container.

\subsection{Inferring contents via remote vibration sensing}
Cameras can sense vibration without contact.
Small motions in video reveal material properties~\cite{bouman2013estimating,davis2015visual} and recover sound~\cite{davis2014visual}, and with a physical model they can map stiffness inside an object~\cite{feng2022visual} or beneath a surface~\cite{ogren2025visual}.
These methods need visible motion and surface texture.
Laser speckle removes both needs.
A defocused camera sees the speckle pattern shift whenever the surface tilts, which makes tiny vibrations easy to track~\cite{zalevsky2009simultaneous}.
Speckle has been used to recover sound~\cite{sheinin2022dual,cai2025event2audio,bagon2026drum} and to locate impacts~\cite{zhang2023analyzing}.
Most relevant to us, Kichler et al.~\cite{kichler2025learning} used a speckle grid to estimate how full an opaque bottle is.
A fill level is one number, and it changes the response of the whole container.
We recover an image of an unknown object instead, whose effect on the wall is local and depends on its shape and position.
To do this, we measure the wall at many points over a wide range of frequencies and train a network, with a camera as the teacher, to turn these measurements into an image.


% =====================================================================
% FLOW 2: narrative, "the wall is the sensor" (through it -> on it -> watching it)
% =====================================================================
\section{Related Work}
\label{sec:related}

There are three broad ways to learn what is inside a closed container: send a wave through its walls, attach sensors that listen to it, or watch it vibrate from a distance.
We review each in turn.

\subsection{Inferring contents via active sensing}
The classic approach treats the container wall as an obstacle and tries to get past it.
X-rays pass through most materials, which is why they are used to screen luggage, but they are ionizing and are stopped by dense metal.
Radio waves are safer.
WiFi can sense people behind a wall~\cite{adib2013see}, and when a camera supplies the training labels, networks learn to predict pose and body shape from radio alone~\cite{zhao2018rfpose,wang2019personinwifi,geng2022densepose}.
Millimeter-wave radar brings this to smaller scales, imaging objects inside envelopes and sealed boxes~\cite{ganse2021envelopes,dodds2025mmnorm,lu2026geraf,dodds2025waveformer}.
Yet radar sees objects the way a flashlight sees a mirror.
Only surfaces that face the antenna send a signal back, so these methods must rebuild the rest through surface-normal estimation~\cite{dodds2025mmnorm}, neural fields~\cite{lu2026geraf}, or learned shape completion~\cite{dodds2025waveformer}.
And if the wall is metal, no radio gets through at all.
Sound~\cite{lindell2019acoustic,christensen2020batvision} and light~\cite{velten2012recovering} can also recover hidden geometry from echoes, but only when there is an open path to the object.

In all of these methods, the wall is what stands in the way.
In ours, the wall is the sensor.
A speaker makes the wall vibrate, and whatever sits inside the container changes that vibration.
We never need to see past the wall, so its material does not matter, and the hardware is just a speaker, a laser, and a camera.

\subsection{Inferring contents via contact vibration sensing}
Listening to vibrations is an old way to learn about objects.
Tap an object and the sound tells you about its shape and material, an idea that has been formalized with physics models~\cite{zhang2017shape} and large multisensory datasets~\cite{gao2023objectfolder,clarke2023realimpact}.
Robots listen through contact microphones to identify what they grasp~\cite{liu2024sonicsense} and where they are touched~\cite{zaielamri2026vibrosense}, and tables and walls become touch sensors the same way~\cite{liu2017vibsense,iravantchi2023sawsense}.
Engineers listen to buildings and plates, finding added masses and damage from shifts in their resonances~\cite{dackermann2013identification,ostachowicz2002location,azuara2021damage}.
In every case the sensor touches the structure, and what comes out is a label, a property, or a point.

The Boombox~\cite{chen2021boombox} is the first to produce an image this way, and it is the work closest to ours.
It records objects falling into a box with contact microphones and learns to draw where they land.
But the Boombox can only hear an object while it is moving.
Once the object settles, the box goes quiet.
We supply the sound ourselves with a speaker, so a resting object is still heard through the way it changes the wall's response, and we listen with a camera instead of a microphone, so nothing touches the container.

\subsection{Inferring contents via remote vibration sensing}
Vibrations can also be seen.
Davis et al.\ showed that a camera can recover sound from the tiny motions of a chip bag~\cite{davis2014visual} and estimate material properties from how objects sway~\cite{davis2015visual}, following earlier work on fabric~\cite{bouman2013estimating}.
With a physical model, such videos can even reveal stiffness inside an object~\cite{feng2022visual} or beneath a surface~\cite{ogren2025visual}.
These methods need motion large enough to see and a surface with texture.

Laser speckle lifts both limits.
A laser spot on a rough surface forms a speckle pattern, and a defocused camera sees that pattern slide whenever the surface tilts, even by a tiny amount~\cite{zalevsky2009simultaneous}.
Speckle has turned everyday objects into microphones~\cite{sheinin2022dual,cai2025event2audio,bagon2026drum} and located impacts on surfaces~\cite{zhang2023analyzing}.
Kichler et al.~\cite{kichler2025learning} took the step closest to ours, aiming a speckle grid at opaque bottles to learn how full they are.
A fill level, however, is a single number that shifts the response of the whole container.
An object sitting somewhere inside changes only part of the response, and the change depends on its shape and position.
We capture this with a grid of laser spots and a broad frequency sweep, and we train a network, with a camera as the teacher, to read an image of the contents out of the wall's vibration.


% =====================================================================
% FLOW 3: requirements (three goals stated up front, each subsection checks them)
% =====================================================================
\section{Related Work}
\label{sec:related}

A practical way to see inside containers should meet three goals.
It should work through any wall, including metal.
It should need neither contact with the container nor motion of the object.
And it should recover what the object looks like and where it is, not just a label.
We review prior work against these goals.

\subsection{Inferring contents via active sensing}
Most methods for seeing inside containers send a wave through the wall.
X-rays are reliable but ionizing, and dense metal blocks them.
Radio waves are safe, and learned models trained with camera labels can recover people's pose and body shape through walls from WiFi~\cite{adib2013see,zhao2018rfpose,wang2019personinwifi,geng2022densepose}.
Millimeter-wave radar can image objects inside envelopes and boxes~\cite{ganse2021envelopes,dodds2025mmnorm,lu2026geraf,dodds2025waveformer}, but most objects reflect radio like a mirror, so only the surfaces facing the antenna are seen and the rest must be inferred~\cite{dodds2025mmnorm,lu2026geraf,dodds2025waveformer}.
Acoustic echoes~\cite{lindell2019acoustic,christensen2020batvision} and optical time-of-flight~\cite{velten2012recovering} recover hidden geometry only when there is an open path to the object.
These methods meet the second and third goals, but not the first, since none of them works through metal.
Our method does, because the signal is the vibration of the wall itself, not a wave that must cross it.

\subsection{Inferring contents via contact vibration sensing}
Vibrations reveal an object's shape and material~\cite{zhang2017shape,gao2023objectfolder,clarke2023realimpact}, what a robot is holding~\cite{liu2024sonicsense}, where a surface was touched~\cite{zaielamri2026vibrosense,liu2017vibsense,iravantchi2023sawsense}, and where a structure carries extra mass or damage~\cite{dackermann2013identification,ostachowicz2002location,azuara2021damage}.
These methods work through any material, but they attach sensors to the structure and output a label, a property, or a location.
The Boombox~\cite{chen2021boombox} reaches the third goal: it predicts an image of objects dropped into a box from contact microphones on its walls.
It still fails the second, because it needs contact, and because its signal comes from the impact, so it cannot see an object at rest.
We excite the container with a speaker and sense the wall optically, which removes both requirements.

\subsection{Inferring contents via remote vibration sensing}
Cameras can measure vibration from a distance, recovering sound~\cite{davis2014visual}, material properties~\cite{bouman2013estimating,davis2015visual}, and stiffness inside objects~\cite{feng2022visual,ogren2025visual} from small motions in video.
Laser speckle makes these measurements far more sensitive, since a defocused camera sees the speckle pattern shift whenever the surface tilts~\cite{zalevsky2009simultaneous}.
Speckle has been used to recover sound~\cite{sheinin2022dual,cai2025event2audio,bagon2026drum} and to locate impacts~\cite{zhang2023analyzing}.
Kichler et al.~\cite{kichler2025learning} meet the first two goals, using a speckle grid to estimate how full an opaque bottle is.
A fill level, however, is a single number that changes the whole container's response, so the third goal remains open.
An object inside a container changes only part of the response, in a way that depends on its shape and position.
By measuring the wall at many points across a wide frequency sweep and training a network with a camera as the teacher, our method recovers an image of the contents and, to our knowledge, is the first to meet all three goals.
